% =============================================================================
% cap02.tex — OBJETIVOS
% =============================================================================

\chapter{OBJETIVOS}\label{cap:objetivos}

O trabalho tem por escopo a verificação estrutural completa de uma edificação
em madeira lamelada colada, desde a caracterização do material até o
detalhamento das ligações e da ancoragem, executada por meio de um programa de
cálculo desenvolvido especificamente para essa finalidade.

\section{OBJETIVO GERAL}\label{sec:objetivo-geral}

Dimensionar e verificar a estrutura de um galpão em madeira lamelada colada com
pórticos de vigas de duplo-bisel, segundo a EN 1995-1-1 \cite{EN1995-1-1:2004},
implementando a cadeia de verificações em um programa de cálculo automatizado e
reprodutível, e apresentando o memorial resultante.

\section{OBJETIVOS ESPECÍFICOS}\label{sec:objetivos-especificos}

Para alcançar o objetivo geral, foram estabelecidos os seguintes objetivos
específicos:

\begin{alineas}
    \item caracterizar o material segundo a EN 14080 \cite{EN14080:2013},
    definindo a classe de resistência, a classe de serviço e os coeficientes
    de modificação e de fluência dela decorrentes;

    \item levantar as ações permanentes e variáveis atuantes na edificação,
    determinando as forças devidas ao vento pela ABNT NBR 6123
    \cite{NBR6123:2023} e
    organizando as combinações de estado-limite último e de serviço conforme
    a EN 1990 \cite{EN1990:2002};

    \item estabelecer a geometria da viga de duplo-bisel por varredura
    paramétrica do espaço de soluções viáveis, identificando a verificação
    que governa cada região desse espaço;

    \item verificar os estados-limites últimos dos elementos de madeira,
    ou seja, viga principal, pilares e terças, quanto à flexão, ao cisalhamento, às
    tensões na cumeeira e na borda inclinada, à estabilidade lateral e à
    flexo-compressão;

    \item verificar o sistema de contraventamento horizontal e vertical e a
    estabilidade global da edificação, incluindo o caminho das forças
    laterais até a fundação;

    \item verificar os estados-limites de serviço, considerando a fluência da
    madeira na classe de serviço adotada;

    \item dimensionar as ligações metálicas e a ancoragem à fundação,
    aplicando critérios distintos de utilização para modos de ruptura frágeis
    e dúcteis;

    \item implementar o conjunto de verificações em programa de código aberto,
    submetido a uma suíte de testes automatizados, capaz de gerar o memorial
    de cálculo a partir dos dados de entrada.
\end{alineas}

O \autoref{qua:objetivos-capitulos} relaciona cada objetivo específico ao
capítulo em que é desenvolvido.

\qquadro
{Correspondência entre os objetivos específicos e os capítulos}
{\footnotesize
\begin{tabular}{|c|p{8.6cm}|c|}\hline
    \textbf{Alínea} & \textbf{Objetivo específico} & \textbf{Desenvolvido em} \\\hline
    a) & Caracterização do material              & \autoref{cap:definicao-problema} \\\hline
    b) & Ações e combinações                     & \autoref{cap:modelagem} \\\hline
    c) & Geometria por varredura paramétrica     & \autoref{cap:predimensionamento} \\\hline
    d) & Estados-limites últimos dos elementos   & \autoref{cap:dimensionamento} \\\hline
    e) & Contraventamento e estabilidade global  & \autoref{cap:dimensionamento} \\\hline
    f) & Estados-limites de serviço              & \autoref{cap:dimensionamento} \\\hline
    g) & Ligações e ancoragem                    & \autoref{cap:detalhamento} \\\hline
    h) & Implementação e testes automatizados    & \autoref{ap:codigo-fonte} \\\hline
\end{tabular}}
{Elaborado pelo autor (2026)}
{objetivos-capitulos}
{A numeração das alíneas corresponde à do \autoref{sec:objetivos-especificos}.}
{}

\section{DELIMITAÇÃO DO TRABALHO}\label{sec:delimitacao}

O trabalho delimita-se à superestrutura em madeira lamelada colada e aos seus
elementos de ligação. A fundação não é dimensionada: o programa entrega os
esforços característicos e de cálculo em cada base, a saber, reação vertical,
força horizontal e esforço de arrancamento, que constituem os dados de entrada
para o projeto geotécnico, executado em etapa posterior e a partir de sondagem
específica do terreno.

Também não integram o escopo o projeto de instalações, o detalhamento de
fechamentos e vedações, o plano de montagem e a análise de custos.

A resistência ao fogo é verificada, porém sob uma limitação que se declara. O
tempo requerido de resistência não é fixado pelo Eurocódigo: decorre do
regulamento nacional de segurança contra incêndio, isto é, da
ABNT NBR 14432 \cite{NBR14432:2001} e da instrução técnica do corpo de bombeiros estadual
aplicável à ocupação. O valor adotado é, portanto, uma declaração de projeto, e
a verificação vale sob essa declaração.

%% TODO: dado ausente — confirmar o tempo requerido de resistência ao fogo
%% junto ao regulamento aplicável à ocupação e à área da edificação. O valor
%% em uso está declarado como parâmetro de origem PROJETO em core/normativo.py.
